% !TeX root=../main.tex

\chapter{جداول کامل نتایج}
\label{app:fulltables}

این پیوست، جداول تکمیلی نتایج فصل چهارم را، برای مراجعهٔ دقیق‌تر به خروجی کامل شبکهٔ آزمایشی، ارائه می‌دهد. \emph{یادآوری:} پیکربندی‌های «بهترین» در جداول زیر، بر مبنای بیشترین بهبود خام \lr{AUROC} تعیین شده‌اند، بدون درنظرگرفتن هزینهٔ آن‌ها از نظر سرگشتگی یا اعتبار خودرمزگذار؛ پیکربندی توصیه‌شدهٔ این پایان‌نامه (بخش~\ref{sec:res-best}) با این پیکربندی‌های «بهترین خام» یکسان نیست.

\section{پیکربندی بهترین و بدترین بهبود خام، به ازای هر حالت انتخاب لایه}
جدول~\ref{tab:app-best-worst} برای هر یک از ده حالت انتخاب لایه، پیکربندی (شرط، $\lambda$) با بیشترین و کمترین میانگین بهبود \lr{AUROC} در لایه(های) هدف را نشان می‌دهد.

\begin{table}[htb]
\centering
\caption{پیکربندی بهترین و بدترین بهبود خام، به ازای هر حالت انتخاب لایه.}
\label{tab:app-best-worst}
% \resizebox{\textwidth}{!}{%
\begin{tabular}{lccrccr}
\toprule
\multirow{2}{*}{لایه(های) هدف} & \multicolumn{3}{c}{بیشترین بهبود خام} & \multicolumn{3}{c}{کمترین بهبود خام} \\
 & شرط & $\lambda$ & بهبود & شرط & $\lambda$ & بهبود \\
\midrule
لایهٔ 0 & \lr{C2} & $0.3$ & $+0.0014$ & \lr{C3} & $1$ & $-0.0096$ \\
لایهٔ 1 & \lr{C3} & $0.1$ & $+0.0012$ & \lr{C3} & $1$ & $-0.0099$ \\
لایهٔ 2 & \lr{C1} & $0$ & $+0.0004$ & \lr{C3} & $1$ & $-0.0347$ \\
لایهٔ 3 & \lr{C4} & $1$ & $+0.0532$ & \lr{C4} & $0.03$ & $-0.0033$ \\
لایهٔ 4 & \lr{C4} & $0.3$ & $+0.0628$ & \lr{C2} & $1$ & $-0.0041$ \\
لایهٔ 5 & \lr{C4} & $0.3$ & $+0.0678$ & \lr{C2} & $0.01$ & $-0.0000$ \\
لایه‌های 2 و 3 & \lr{C2} & $1$ & $+0.0104$ & \lr{C2} & $0.1$ & $-0.0142$ \\
لایه‌های 3 و 4 & \lr{C3} & $1$ & $+0.0445$ & \lr{C2} & $0.01$ & $-0.0030$ \\
لایه‌های 4 و 5 & \lr{C4} & $0.3$ & $+0.0572$ & \lr{C2} & $0.01$ & $-0.0053$ \\
لایه‌های 3, 4 و 5 & \lr{C3} & $1$ & $+0.0694$ & \lr{C2} & $0.01$ & $-0.0080$ \\
\bottomrule
\end{tabular}%
% }
\end{table}

نکتهٔ قابل‌توجه در این جدول، حضور مکرر شرط
\lr{C2} (فضای خام) در ستون «کمترین بهبود» با مقادیر
$\lambda$
پایین است؛ این حضور، نشانهٔ خطا یا رفتار آسیب‌شناختی نیست، بلکه بازتاب همان یافتهٔ بخش~\ref{sec:res-condition} است که
\lr{C2}
به‌طور میانگین نزدیک به صفر عمل می‌کند و بنابراین، در نبود یک اثر نظام‌مند قوی، کدام پیکربندی خاص کمترین بهبود را داشته باشد را عمدتاً نوسان آماری کوچک تعیین می‌کند، نه یک سازوکار زیان سیستماتیک (برخلاف لایهٔ دوم در شرایط \lr{C3}/\lr{C4} که زیانش بزرگ، پیوسته، و در بخش~\ref{sec:res-layer} توضیح داده شد). توجه شود که پیکربندی توصیه‌شدهٔ این پایان‌نامه
(\lr{C4}، لایهٔ ۵، $\lambda=0.3$) در سطر «لایهٔ ۵» همان پیکربندی بیشترین بهبود خام نیز هست؛ در بقیهٔ سطرها، پیکربندی بیشترین بهبود خام لزوماً با معیارهای بخش~\ref{sec:res-best} سازگار نیست.

\section{میانگین بهبود بر حسب شرط و $\lambda$}
جدول~\ref{tab:app-lambda} میانگین بهبود \lr{AUROC} (در لایه(های) هدف، میانگین‌گیری‌شده روی هر ده حالت انتخاب لایه) را برای هر ترکیب شرط و $\lambda$ نشان می‌دهد؛ این جدول، دادگان زیرین شکل~\ref{fig:gain-vs-lambda} است.

\begin{table}[htb]
\centering
\caption{میانگین بهبود \lr{AUROC}، به تفکیک شرط و $\lambda$.}
\label{tab:app-lambda}
\begin{tabular}{lccc}
\toprule
$\lambda$ & \lr{C2} (خام) & \lr{C3} (\lr{SAE}) & \lr{C4} (\lr{SAE}+حمایت) \\
\midrule
$0.01$ & $-0.0036$ & $-0.0019$ & $-0.0014$ \\
$0.03$ & $-0.0028$ & $+0.0063$ & $+0.0053$ \\
$0.10$ & $-0.0034$ & $+0.0267$ & $+0.0237$ \\
$0.30$ & $-0.0007$ & $+0.0248$ & $+0.0341$ \\
$1.00$ & $+0.0043$ & $+0.0300$ & $+0.0291$ \\
\bottomrule
\end{tabular}
\end{table}

برای مقایسه،
\lr{C1} (بدون $\lambda$، اجرای منفرد) میانگین بهبود $-0.0025$ را در همین لایه‌ها به دست می‌آورد.

\section{اعتبار بازسازی خودرمزگذار در گام پایانی}
جدول~\ref{tab:app-fidelity} میانگین واریانس تبیین‌شدهٔ بازسازی را در گام پایانی آموزش، به تفکیک شرط و
$\lambda$، تکمیل‌کنندهٔ شکل~\ref{fig:fidelity}، نشان می‌دهد.

\begin{table}[htb]
\centering
\caption{میانگین واریانس تبیین‌شدهٔ بازسازی در گام پایانی، به تفکیک شرط و $\lambda$.}
\label{tab:app-fidelity}
\begin{tabular}{lccc}
\toprule
$\lambda$ & \lr{C2} (خام) & \lr{C3} (\lr{SAE}) & \lr{C4} (\lr{SAE}+حمایت) \\
\midrule
$0.01$ & $0.9629$ & $0.9613$ & $0.9618$ \\
$0.03$ & $0.9631$ & $0.9503$ & $0.9529$ \\
$0.10$ & $0.9625$ & $0.9304$ & $0.9302$ \\
$0.30$ & $0.9594$ & $0.9200$ & $0.9225$ \\
$1.00$ & $0.9163$ & $0.8929$ & $0.9091$ \\
\bottomrule
\end{tabular}
\end{table}

برای مقایسه، شرط
\lr{C1}
(فاقد هرگونه هدف ناوردایی) میانگین $0.9634$ را به دست می‌آورد که نزدیک‌ترین مقدار به بالاترین ردیف جدول بالا (کمترین $\lambda$) است؛ این نزدیکی، خود، شاهدی است بر اینکه در
$\lambda$
بسیار کوچک، رفتار آموزش عملاً به رفتار شرط کنترلی همگرا می‌شود. سطر
$\lambda=0.30$
در این جدول--پیکربندی توصیه‌شدهٔ این پایان‌نامه (بخش~\ref{sec:res-best}) دقیقاً از همین سطر، ستون
\lr{C4}، برآمده است--نشان می‌دهد در این نقطه، واریانس تبیین‌شده هنوز به‌خوبی بالاتر از آستانهٔ هشدار
$90\%$
باقی می‌ماند؛ برخلاف سطر
$\lambda=1.00$
که ستون
\lr{C3}
آن، از این آستانه عبور می‌کند.
